% =====================================================================
% New Appendix B subsection: proofs of Lemma 1, Proposition 1, Proposition 2
% Insert as the first subsection of Appendix B (before "A. Conditional
% Objectives and Update Model"), then in the main text replace each
% "given in the supplementary material [[TODO ...]]" with
% "given in Appendix~\ref{app:mech_proofs}".
% =====================================================================

\subsection{Proofs of the Mechanism-Level Results}
\label{app:mech_proofs}

All squares, square roots, and absolute values of vectors below are taken
coordinate-wise.

\paragraph*{Proof of Lemma~1}
Let $\mathcal{I}_\ell$ be the feature set of source $\ell$, with
$|\mathcal{I}_\ell| = n_\ell$, and let $\mathcal{I} = \bigcup_\ell \mathcal{I}_\ell$,
$|\mathcal{I}| = n$. Because the sets are disjoint,
\begin{equation*}
\mu_{\mathrm{pool}} = \frac{1}{n}\sum_{\ell=1}^{S}\sum_{i\in\mathcal{I}_\ell} z_i
= \sum_{\ell=1}^{S}\frac{n_\ell}{n}\,\mu_\ell = \sum_{\ell=1}^{S} w_\ell\,\mu_\ell .
\end{equation*}
For the variance, decompose each deviation as
$z_i - \mu_{\mathrm{pool}} = (z_i - \mu_\ell) + (\mu_\ell - \mu_{\mathrm{pool}})$ for
$i \in \mathcal{I}_\ell$. Squaring and summing over $\mathcal{I}_\ell$, the
cross term vanishes because $\sum_{i\in\mathcal{I}_\ell}(z_i - \mu_\ell) = 0$.
With the population convention in~(1),
\begin{equation*}
\sum_{i\in\mathcal{I}_\ell}(z_i - \mu_{\mathrm{pool}})^2
= n_\ell\,\sigma_\ell^2 + n_\ell\,(\mu_\ell - \mu_{\mathrm{pool}})^2 .
\end{equation*}
Summing over $\ell$ and dividing by $n$ gives
$\sigma_{\mathrm{pool}}^2 = \sum_\ell w_\ell\bigl[\sigma_\ell^2 + (\mu_\ell - \mu_{\mathrm{pool}})^2\bigr]$,
which is~(2). \hfill$\blacksquare$

\paragraph*{Proof of Proposition~1}
For diagonal Gaussians, (8) gives
$W_2^2(\widetilde{G}, G_*) = \|\tilde\mu - \mu_*\|_2^2 + \|\tilde\sigma - \sigma_*\|_2^2$.
Write $\Delta\mu = \mu_* - \mu_m$.

\emph{Mean term.} Since $\mu_c = \mu_*$, (15) gives
$\tilde\mu - \mu_* = (1-\alpha)(\mu_m - \mu_*)$, hence
\begin{equation*}
\|\tilde\mu - \mu_*\|_2^2 = (1-\alpha)^2\|\Delta\mu\|_2^2 .
\end{equation*}

\emph{Spread term.} For $a, b \ge 0$,
$(a-b)^2 \le |a-b|\,(a+b) = |a^2 - b^2|$. Applying this to each coordinate of
$\tilde\sigma$ and $\sigma_*$,
\begin{equation*}
\|\tilde\sigma - \sigma_*\|_2^2 \le \bigl\|\tilde\sigma^2 - \sigma_*^2\bigr\|_1 .
\end{equation*}
With $\sigma_c = \sigma_*$, (15) gives
\begin{equation*}
\tilde\sigma^2 - \sigma_*^2
= (1-\alpha)(\sigma_m^2 - \sigma_*^2) + \alpha(1-\alpha)\,\Delta\mu^2 + \zeta\mathbf{1},
\end{equation*}
so by the triangle inequality,
\begin{equation*}
\bigl\|\tilde\sigma^2 - \sigma_*^2\bigr\|_1
\le (1-\alpha)\,\|\sigma_m^2 - \sigma_*^2\|_1 + \alpha(1-\alpha)\,\|\Delta\mu\|_2^2 + d\zeta .
\end{equation*}

\emph{Combination.} Adding the two terms and using
$(1-\alpha)^2 + \alpha(1-\alpha) = 1-\alpha$,
\begin{equation*}
W_2^2(\widetilde{G}, G_*) \le (1-\alpha)\bigl(\|\mu_* - \mu_m\|_2^2 + \|\sigma_m^2 - \sigma_*^2\|_1\bigr) + d\zeta,
\end{equation*}
which is~(11). \hfill$\blacksquare$

\paragraph*{Proof of Proposition~2}
Let $M_\ell \in \{0,1\}$ indicate that the mask initially marks output $\ell$
for removal; the $M_\ell$ are independent with $\Pr(M_\ell = 1) = p$. Let
$\mathcal{S}_j$ be the $S_j \ge 1$ outputs that contain class~$j$.

\emph{Class loss.} Class $j$ is absent after masking if and only if every
output in $\mathcal{S}_j$ is finally removed. This requires the event
$A_j = \{M_\ell = 1\ \forall \ell \in \mathcal{S}_j\}$, with
$\Pr(A_j) = p^{S_j}$, and excludes the case in which all $S$ outputs are
marked and the uniformly restored output belongs to $\mathcal{S}_j$. That case
has probability $p^{S}\cdot S_j/S$ and is contained in $A_j$. Therefore
\begin{equation*}
\Pr\bigl[j \notin \mathcal{J}_t\ \text{after masking}\bigr] = p^{S_j} - \frac{S_j}{S}\,p^{S},
\end{equation*}
which is~(12). For $S = 1$ the expression equals zero, consistent with the
deterministic retention of a single output.

\emph{Expected retained outputs.} Each output survives the initial marking
with probability $1-p$, giving $S(1-p)$ in expectation. Restoration adds
exactly one output on the event that all $S$ outputs are marked, which has
probability $p^S$. Hence the expected number of retained outputs is
$S(1-p) + p^S$. \hfill$\blacksquare$
